\chapter{Conclusioni}
In questo capitolo si riportano le conclusioni tratte da quanto discusso nei capitoli precedenti
di questo elaborato sui temi principali affrontati, ovvero:
\begin{itemize}
    \item le implicazioni organizzative della Direttiva NIS2;
    \item il ruolo dei principali framework e standard internazionali nell'adozione di tale direttiva;
    \item il ruolo della figura del CISO all'interno dell'azienda, le sue responsabilità e i suoi compiti
        in relazione all'adozione di NIS2.
\end{itemize}
Inoltre, essendo la Direttiva NIS2 di recente creazione, si aggiunge una breve riflessione sul futuro impatto
di tale normativa sullo scenario europeo e sul ruolo delle nuove evoluzioni tecnologiche, come l'intelligenza
artificiale, nel contesto della cybersicurezza.



% 6.1 %%%%%%%%%%%%%%%%%%%%%%%%%%%%%%%%%%%%%%%%%%%%%%%%%%%%%%%%%%%%%%%%%%%%%%%%
\section{Implicazioni organizzative della Direttiva NIS2}
\subsection{Responsabilità nella gestione del rischio}
La principale novità introdotta da NIS2 è certamente l'attribuzione della responsabilità
della sicurezza informatica agli organi di gestione dell'azienda. Infatti, pur essendo già
responsabile di tutto ciò che avviene all'interno dell'organizzazione, con questa misura si
formalizza il ruolo del CdA come supervisore delle misure di cybersicurezza e la non delegabilità
di tale funzione a organi aziendali sottoposti. Inoltre, i componenti dell'amministrazione sono 
anche direttamente sanzionabili in caso di inadempimenti normativi in materia di NIS2, con cospicue 
ammende economiche e sospensioni dal proprio incarico nelle casistiche di maggiore gravità.

Le implicazioni di questo provvedimento sono molteplici e si riflettono direttamente sulla funzione
di cybersicurezza aziendale e, di conseguenza, sul CISO. Tale figura si trova dunque formalmente 
sollevata da responsabilità dirette riguardanti l'inadempimento normativo, ma d'altro canto deve
assicurare la corretta implementazione delle misure richieste, individuare eventuali lacune e
coadiuvare il CdA nella conformità normativa. La gestione della cybersecurity
da parte dell'amministrazione richiede infatti una costante informazione di tali organi, implicando la 
necessità del CISO di interagire trasversalmente con le altre funzioni aziendali per individuare e 
valutare le vulnerabilità informatiche. A tal fine, le aziende devono implementare nuovi framework
di gestione del rischio che contemplino la sicurezza informatica come componente centrale, nonché 
specifici meccanismi di periodico dialogo e formazione in materia.


\subsection{Misure di sicurezza e sicurezza della supply chain}
Un ulteriore punto cardine della Direttiva NIS2 è rappresentato dalla gestione della sicurezza in 
maniera integrata lungo tutta la filiera produttiva di ogni azienda. Tale disposizione mira a garantire un livello uniforme di sicurezza tra tutte le aziende, assicurando in particolare
maggiori livelli di protezione e resilienza verso minacce legate alle interazioni con l'esterno.

NIS2 richiede ai fornitori dei soggetti che rientrano nel campo di applicazione della normativa di implementare 
meccanismi analoghi a quelli dei clienti per garantire sicurezza lungo tutta la filiera. Ciò si applica, secondo
il legislatore, non solo ai fornitori diretti ma anche a tutte le entità identificate come critiche all'interno
della catena. 

Concretamente, la misura è implementata dalle organizzazioni mediante clausole apposite inserite nel contratto, 
che assicurino il \textit{commitment} del fornitore per garantire un livello di cybersicurezza conforme. Inoltre,
i vincoli contrattuali devono essere affiancati da meccanismi di controllo applicati direttamente dai clienti
ai propri partner critici, quali richieste di certificazioni e auditing, valutando periodicamente il profilo
di rischio del soggetto.

Il coinvolgimento attivo dei fornitori nel controllo della cybersicurezza estende notevolmente
il raggio d'azione della normativa, già superiore rispetto alla precedente NIS. Se da un lato questo è desiderabile
per il legislatore, che mira alla massima diffusione di un livello di sicurezza uniforme, dall'altro rappresenta
un ingente aumento di costi per imprese che in precedenza non disponevano di meccanismi di cybersecurity adeguati,
mentre ora necessitano di sviluppare metodologie e framework appositi coerentemente con quanto richiesto da NIS2.


\subsection{Notifica degli incidenti}
NIS2 introduce anche meccanismi dettagliati per la notifica di eventi di sicurezza informatica alle autorità. In questo 
modo il legislatore vuole favorire il monitoraggio della sicurezza da parte degli organi nazionali preposti a tale
funzione e, al tempo stesso, fornire alle aziende una metodologia strutturata per la gestione degli incidenti. Le 
agenzie governative svolgono quindi un ruolo fondamentale per la raccolta delle informazioni su minacce e attacchi
informatici e per l'assistenza e il coordinamento dei soggetti durante le crisi.

La normativa fissa precisi indicatori chiave di prestazione (KPI) per la segnalazione di incidenti significativi, stabilendo periodici aggiornamenti 
dopo ventiquattro e settantadue ore dal verificarsi dell'evento. Trascorso un mese, è richiesto alle imprese un 
rapporto finale sugli avvenimenti, promuovendo così l'apprendimento e la collaborazione tra le varie aziende, 
individuando e catalogando le minacce e i relativi piani di risposta.

Per implementare queste disposizioni le imprese devono includere nei propri piani di gestione delle crisi il 
costante dialogo con le autorità. A tale scopo, la normativa richiede di definire un punto di contatto interno
alle organizzazioni che si occupi di dialogare con gli organi governativi e le entità esterne in caso di incidenti,
funzione in capo al CISO. 


\subsection{Sfide per le organizzazioni}
Le aree su cui le organizzazioni devono agire per adempiere completamente a NIS2 sono molteplici e trasversali
rispetto ai differenti core business esistenti e richiedono quindi investimenti nella gestione della cybersicurezza
in modo integrato a livello aziendale.

Appare chiaro come, in questo contesto, le entità di dimensioni maggiori o già oggetto di precedenti normative di 
sicurezza (quali NIS e DORA) siano avvantaggiate in questo processo: esse dispongono di risorse maggiori e spesso
sono già dotate di framework di cybersecurity interni. Questi soggetti devono dunque limitarsi a condurre una 
\textit{gap analysis} tra le misure implementate e quelle richieste dal legislatore europeo, verificando il 
proprio profilo e predisponendo investimenti mirati per consolidarlo.

I soggetti più colpiti dal punto di vista delle risorse sono quindi le entità non oggetto di precedenti normative
o di dimensioni inferiori, con capacità finanziarie inferiori e minore \textit{know-how} in materia. Esse devono
infatti costruire il proprio framework \textit{ex novo}, assumendo figure esperte e modificando i propri processi
in modo da soddisfare gli obblighi di NIS2. Ciò comporta ingenti investimenti che, tuttavia, non sempre sono agevoli
per aziende in questa situazione. Standard come ISO 27001 e NIST CSF possono rappresentare validi strumenti su cui
sviluppare i propri meccanismi di gestione della cybersicurezza, da integrare adeguatamente rispetto a NIS2, 
rendendo il processo più veloce e contenendo i costi.

\begin{table}[H]
    \centering
    % riduco leggermente spazi e aumento interlinea della tabella
    % (richiede xcolor con opzione table: \usepackage[table]{xcolor})
    \setlength{\tabcolsep}{5pt}
    \renewcommand\arraystretch{1.1}
    \begin{tabularx}{\textwidth}{|>{\raggedright\arraybackslash}p{0.15\textwidth}|>{\raggedright\arraybackslash}p{0.17\textwidth}|>{\raggedright\arraybackslash}p{0.15\textwidth}|>{\raggedright\arraybackslash}p{0.13\textwidth}|>{\raggedright\arraybackslash}X|}
    \hline
        \rowcolor{gray!20}
        \textbf{Dimensione aziendale} & \textbf{Disponibilità di risorse} & \textbf{Maturità dei sistemi} & \textbf{Impatto NIS2} & \textbf{Note} \\ \hline
        Piccola & Bassa & Assente o minima & \cellcolor{red!25}{Molto elevato} & \small{Assenza di funzioni dedicate, carenza di policy e controlli, mancanza di personale qualificato.} \\ \hline
        Piccola-media & Media & Iniziale & \cellcolor{red!15}{Elevato} & \small{Alcune misure implementate, ma governance e risk management deboli.} \\ \hline
        Media & Media & Intermedia & \cellcolor{yellow!20}{Medio} & \small{Cybersicurezza strutturata, ma approccio parziale e frammentato.} \\ \hline
        Medio-grande & Elevata & Elevata & \cellcolor{green!15}{Basso} & \small{Gestione del rischio formalizzata, framework interni implementati, necessità di adeguamento agli obblighi di notifica.} \\ \hline
        Grande & Elevata & Completa & \cellcolor{green!25}{Molto basso} & \small{Già oggetto di regolamentazione, compliance e interazioni con le autorità integrati nei processi, gap analysis per eventuali carenze.} \\ \hline
    \end{tabularx}
    \vspace{0.6\baselineskip}
    \caption{Impatto atteso della Direttiva NIS2 per classi dimensionali e maturità organizzativa}
    \label{tab:impatto-nis2}
\end{table}

La tabella \ref{tab:impatto-nis2} espone le considerazioni precedenti in una visione sintetica, delineando l'impatto della 
Direttiva NIS2 sulle aziende in base alla loro dimensione, alle risorse di cui dispongono e alla maturità dei sistemi di
sicurezza informatica precedentemente adottati.



% 6.2 %%%%%%%%%%%%%%%%%%%%%%%%%%%%%%%%%%%%%%%%%%%%%%%%%%%%%%%%%%%%%%%%%%%%%%%%
\section{Ruolo degli standard internazionali nell'adempimento della Direttiva NIS2}
\subsection{ISO/IEC 27001:2022 a confronto con NIS2}
La certificazione ISO 27001 rappresenta una delle più adottate a livello mondiale nell'ambito della cybersicurezza.
Lo standard fornisce alle aziende una struttura di governance solida e un insieme di controlli di sicurezza da adottare 
per gestire e mitigare il rischio informatico.

Il principale vantaggio che deriva dall'utilizzo di ISO 27001 relativamente all'adempimento di NIS2 è la governance.
ISO fornisce infatti una governance profondamente strutturata e una documentazione approfondita dei processi e delle misure
attuate per contrastare le minacce informatiche, imponendo anche una struttura per la definizione delle responsabilità.
La certificazione ISO rappresenta quindi per le organizzazioni uno strumento cruciale di accountability che, con un costo 
relativamente contenuto, permette di dimostrare l'impegno nella cybersicurezza a fronte di audit da parte di autorità esterne.

Anche i controlli dell'\textit{Annex A} coprono le aree individuate come critiche da NIS2, costituendo quindi una solida 
base per soddisfare gli obblighi normativi: tra essi si ritrovano misure in materia di formazione del personale, \textit{threat intelligence},
continuità operativa e ripristino in caso di incidenti.

ISO è dunque coerente con buona parte della normativa europea, ma non ne soddisfa completamente i requisiti richiesti. Infatti, sebbene 
alcuni framework di cybersicurezza nazionali equiparino l'adempimento ISO a quello NIS2, i due strumenti differiscono sotto più punti di vista.
Si ricorda, inolre, la possibilità di colmare alcune lacune dello standard 27001 rispetto a NIS2 integrando altri standard ISO.

In primo luogo, la certificazione ISO 27001 attesta solamente la presenza di adeguati accorgimenti di sicurezza informatica, ma non ne garantisce
l'effettiva efficacia: la modalità implementativa del controllo è lasciata alla singola azienda, la quale deve definirla sulla base del profilo di 
rischio da essa desiderato. NIS2, al contrario, richiede un'implementazione delle misure effettiva e basata sul livello di rischio ricoperto 
dall'impresa all'interno del sistema Paese e della propria catena produttiva. Dunque, la certificazione ISO permette alle aziende di dimostrare
l'adozione di determinati mezzi di contrasto alle minacce informatiche, ma l'effettiva riduzione del rischio apportata da essi deve essere 
verificata in relazione al profilo individuato da NIS2 per l'azienda.

Inoltre, data la natura non regolatoria dello standard, ISO non prevede gli obblighi di notifica presenti invece nella direttiva, così
come l'attribuzione delle responsabilità e le relative sanzioni verso gli organi di gestione non trovano diretta corrispondenza.

Infine, sebbene ISO presenti controlli dedicati a quest'area di rischio, le misure di sicurezza della catena del valore previste
da NIS2 risultano più ampie e pervasive, riguardando anche la necessità di controllo diretto da parte del soggetto sui propri fornitori.


\subsection{NIST CSF 2.0 a confronto con NIS2}
Il framework NIST e la Direttiva NIS2 condividono una struttura simile. Il NIST CSF 2.0 ha infatti una vocazione più operativa rispetto a 
ISO 27001, rendendolo più aderente alle richieste del legislatore europeo. 

Il NIST CSF si presenta come un framework ciclico, ma in esso si delineano due parti ben definite:
\begin{itemize}
    \item le funzioni \textit{Identify} e \textit{Protect} delineano la normale condizione di operatività dell'impresa e le misure di base
        per la sicurezza informatica;
    \item le funzioni \textit{Detect}, \textit{Respond} e \textit{Recover} trattano invece il modo in cui l'organizzazione identifica 
        e reagisce agli attacchi cui è soggetta.
\end{itemize}
Tutte le funzioni si ergono sulla funzione \textit{Govern}, ponendo la governance al centro della gestione della cybersicurezza aziendale.

L'assetto del framework NIST si riflette nelle misure richieste da NIS2 che, a loro volta, mirano ad abbracciare completamente
tutti gli aspetti di cybersicurezza che interessano l'organizzazione, in ogni fase della sua operatività. Inoltre, gli obiettivi stabiliti
dal NIST all'interno delle singole funzioni trovano piena corrispondenza tra le misure richieste dal legislatore europeo, così come la
gestione dei rischi informatici lungo la filiera produttiva, seppur senza obblighi di controllo espliciti sui fornitori.

Nonostante i numerosi punti di contatto, NIS2 e NIST CSF presentano anche alcune divergenze significative. 

Nuovamente, trattandosi di un framework volontario e non di uno strumento di natura regolatoria, gli obblighi di notifica alle autorità
e la comunicazione con entità esterne per la gestione degli eventi di cybersicurezza
non sono pienamente contemplati nel framework NIST. Esso, infatti, si limita a segnalare la necessità di adottare meccanismi di avviso 
ai soggetti interessati da eventuali incidenti, mentre NIS2 specifica indicatori espliciti (KPI) e richiede l'implementazione di un 
piano di gestione delle emergenze che coinvolga le autorità ed eventuali terzi danneggiati dagli eventi.

Tuttavia, il limite più grande del NIST CSF è rappresentato dalla mancanza di una certificazione riconosciuta delle misure adottate 
conformemente al framework: ciò implica la perdita della possibilità di dimostrare accountability per l'organizzazione, costituendo quindi
un punto di profonda distanza rispetto agli obblighi NIS2 per la rendicontazione dei controlli implementati in caso di audit da parte delle 
agenzie governative di controllo.


\subsection{Ruolo degli standard internazionali nell'adempimento di NIS2}
La Direttiva NIS2 rappresenta un primo passo del legislatore nella direzione di un framework europeo unico per la cybersicurezza. Infatti, 
seppur mantenga un carattere normativo, delinea le aree in cui agire per assicurare la sicurezza informatica dell'organizzazione,
sia per quanto riguarda i controlli sia per l'assegnazione delle responsabilità.

Da questo punto di vista, dunque, un'azienda che si trova per la prima volta a dover implementare misure di sicurezza per adempiere agli 
obblighi NIS2 può decidere di basarsi direttamente sulla direttiva, sul relativo regolamento operativo e sulle norme di recepimento
nazionale -- che descrivono i controlli nel dettaglio -- per realizzare il proprio framework di cybersicurezza, senza la necessità di 
adottare standard intermedi. Se realizzato correttamente, ciò favorisce il pieno adempimento normativo sin da subito, senza il bisogno di condurre
analisi per l'identificazione di eventuali inadempienze; inoltre, l'azienda potrà decidere anche in seguito di adeguare le misure implementate
per aderire a uno o più standard.

La figura Fig. \ref{fig:impl-nis2-diretta}, basata sul diagramma proposto in \cite{parmar-miles-csf}, riassume in un diagramma di flusso il processo 
di implementazione della direttiva NIS2 senza l'ausilio di framework.

\defaultimage{impl-nis2-diretta}{Flusso di implementazione NIS2 senza l'utilizzo di framework.}{fig:impl-nis2-diretta}

Tuttavia, è importante notare che questo è un processo particolarmente complesso e oneroso, soprattutto per le imprese che non hanno 
un \textit{know-how} interno in merito ai rischi informatici e necessitano quindi di assumere figure esperte e costruire un proprio sistema
senza una base preesistente. In questo contesto, adottare lo standard ISO 27001 o il framework NIST CSF 2.0 può rappresentare un'opzione valida, in quanto:
\begin{itemize}
    \item rende il processo di adeguamento più rapido, in particolare nella parte iniziale, per poi portare a una \textit{gap analysis} rispetto agli obblighi NIS2;
    \item favorisce lo sviluppo del framework di cybersicurezza interno all'azienda, basato su pratiche consolidate nelle altre organizzazioni.
\end{itemize}

La figura Fig. \ref{fig:impl-nis2-framework}, basata sul diagramma proposto in \cite{parmar-miles-csf}, mostra come un'azienda può raggiungere 
l'adempimento di NIS2 tramite l'utilizzo degli standard citati, integrandoli nella propria strategia di cybersicurezza e adottando misure correttive 
per gli obblighi normativi non coperti.

\defaultimage{impl-nis2-framework}{Flusso di implementazione NIS2 con l'utilizzo di framework.}{fig:impl-nis2-framework}

In sostanza, l'adozione diretta della direttiva NIS2, piuttosto che la realizzazione di un framework iniziale basato su standard internazionali largamente
diffusi, dipende principalmente dalla maturità dei processi di gestione del rischio informatico in essere.

Aziende che dispongono di meccanismi di sicurezza consolidati e che, con ogni probabilità, già adottano ISO o NIST troveranno più naturale condurre
un'analisi che attesti le differenze tra il proprio profilo e quello delineato dalla Direttiva NIS2, realizzando quindi investimenti mirati 
nelle aree individuate come critiche e rinforzando quanto già realizzato.

Al contrario, le organizzazioni che devono costruire la propria gestione della cybersecurity devono valutare, in base alle proprie risorse, se risulti
più conveniente realizzarla direttamente sulla base del materiale normativo o se sia preferibile una fase di transizione intermedia che coinvolga
l'adozione di uno standard di sicurezza già consolidato, per poi colmare gli inadempimenti.



% 6.3 %%%%%%%%%%%%%%%%%%%%%%%%%%%%%%%%%%%%%%%%%%%%%%%%%%%%%%%%%%%%%%%%%%%%%%%%
\section{Ruolo del CISO nell'adempimento della Direttiva NIS2}
\subsection{Evoluzione della figura del CISO}
La figura del \textit{Chief Information Security Officer} in azienda è impegnata principalmente su cinque fronti:
\begin{enumerate}
    \item \textbf{Governance e gestione del rischio}
    \item \textbf{Gestione della security by design}
    \item \textbf{Gestione delle operations}
    \item \textbf{Gestione dell'innovazione}
    \item \textbf{Training e awareness}
\end{enumerate}
Il CISO ricopre quindi un ruolo estremamente trasversale rispetto a tutte le funzioni aziendali e le sue mansioni sono profondamente
cambiate nel recente periodo, in seguito all'introduzione delle moderne tecnologie e ai tentativi del legislatore di regolarne l'utilizzo
per tutelare le organizzazioni e i consumatori. Il CISO è passato quindi da svolgere compiti di carattere prettamente tecnico alla necessità
di integrare nel proprio profilo competenze di leadership e di gestione del rischio, coordinando le proprie azioni con il piano di business aziendale.

In questo contesto si colloca la necessità di assicurare la conformità alla Direttiva NIS2, con un impatto diretto su tutto il team di cybersecurity aziendale.
Sebbene non vi sia un aumento delle responsabilità dirette del CISO, egli si trova a dover interagire frequentemente con gli organi di gestione, per adempiere
agli obblighi formativi e informativi, e con i referenti delle altre funzioni, per individuare e gestire eventuali vulnerabilità.


\subsection{Responsabilità e gestione del rischio}
La Direttiva NIS2 sancisce la responsabilità diretta degli organi di gestione in materia di sicurezza informatica aziendale, mentre figure tecniche come il CISO
non sono imputabili in caso di inadempimenti normativi e non possono essere oggetto di deleghe da parte del CdA su queste tematiche.
In materia NIS2, il CISO si occupa di individuare le potenziali vulnerabilità, valutare i rischi in relazione al risk framework aziendale, gestirli in 
cooperazione con il rispettivo risk owner e sovrintendere l'implementazione della formazione, delle policy e delle misure di sicurezza informatica necessarie
per la conformità agli obblighi normativi. 

Pur essendo vero che NIS2 non introduce nessuna nuova responsabilità, la direttiva implica per il CISO nuovi obblighi sul piano informativo e relazionale:
è necessario informare periodicamente il CdA sul profilo di rischio dell'azienda in materia di sicurezza informatica, dialogare con i responsabili delle 
funzioni aziendali per individuare e gestire le fonti di potenziali attacchi e strutturare un piano di formazione adeguato per ogni livello
della gerarchia organizzativa. La sicurezza rappresenta quindi un tema non solo tecnico, ma soprattutto strategico e di governance organizzativa, richiedendo 
il coinvolgimento attivo del team di cybersicurezza nella definizione della strategia delle varie aree di business e nel dialogo periodico con il CdA e 
il top management.

In quest'ottica, gran parte del lavoro svolto dal CISO oggi riguarda il \textit{risk management}: il compito principale dell'area di cybersecurity rimane
l'individuazione delle vulnerabilità e la gestione dei rischi correlati. In questo processo, è necessario valutare le potenziali fonti di minacce in relazione
al risk framework dell'organizzazione e al livello di rischio considerato accettabile dal legislatore, collaborando con il rispettivo \textit{risk owner} per
trattare o accettare i rischi non eliminabili direttamente.



% 6.4 %%%%%%%%%%%%%%%%%%%%%%%%%%%%%%%%%%%%%%%%%%%%%%%%%%%%%%%%%%%%%%%%%%%%%%%%
\section{Considerazioni generali e visione prospettica}
\subsection{Impatto futuro di NIS2}
La Direttiva NIS2, soprattutto a causa del grande numero di soggetti che coinvolge, sta avendo un impatto su larghissima scala, coinvolgendo 
organizzazioni profondamente diverse per \textit{core business} e dimensioni. Gli effetti di tale normativa sono ancor più evidenti per le 
realtà più piccole e i settori non precedentemente regolamentati, i quali vedranno un crescente aumento degli investimenti in sistemi di 
sicurezza.

Ciò è certamente un fattore positivo, che si allinea con l'obiettivo del legislatore di prioritizzare la sicurezza informatica e diffondere
pratiche, controlli e principi comuni in ogni campo. In questo modo, le aziende sono costrette a trattare la cybersicurezza come un tema 
centrale nella propria strategia, mentre il cospicuo aumento degli investimenti nei sistemi di sicurezza può, almeno sulla carta,
favorirne la continua evoluzione.

Come detto, questa normativa ha anche implicazioni negative, richiedendo ingenti risorse per l'implementazione di meccanismi che soddisfino i 
requisiti richiesti. Tuttavia, ciò appare giustificato se osservato nell'ottica di garantire la sicurezza della società intera: la vera innovazione 
introdotta dalla Direttiva NIS2 risiede nel riconoscimento della sicurezza informatica come responsabilità collettiva, che coinvolge l’intera collettività, 
dalle imprese, chiamate a implementare adeguate misure di protezione, fino ai cittadini, sempre più esposti alle conseguenze degli attacchi digitali.


\subsection{Intelligenza artificiale e sicurezza informatica}
L’intelligenza artificiale è un elemento di trasformazione profonda per la società moderna, e il suo ruolo nella sicurezza informatica è 
una questione aperta e in continua evoluzione.

Come evidenziato dai report ENISA e Clusit e rimarcato anche dai CISO intervistati, l'intelligenza artificiale non ha ancora rivoluzionato 
il campo delle minacce cyber, ma è ragionevole ritenere che sia questione di tempo prima che ciò accada. Nell'ultimo periodo, infatti, l'AI 
è stata largamente impiegata negli attacchi informatici, soprattutto nelle aree del phishing e del social engineering, generando un notevole 
aumento del volume degli attacchi, ma non apportando alcun miglioramento qualitativo considerevole. Dunque, i cybercriminali riconoscono l'importanza 
dell'AI per i loro obiettivi, ma quanto osservato porta a ritenere che ancora oggi le tecnologie non siano sufficientemente sviluppate per 
poter attaccare in maniera completamente autonoma sistemi di difesa complessi.

Analogamente, l'intelligenza di carattere difensivo è un tema fonte di grande discussione: molte aziende intendono puntare sull'AI per costruire
i propri sistemi difensivi del futuro, ma sono pochi gli esempi in cui ciò avviene realmente, mentre sono più i casi 
in cui il termine AI viene utilizzato impropriamente in relazione ai sistemi di cybersicurezza. Investire in queste soluzioni è particolarmente
oneroso e richiede risorse che non sono appannaggio della maggior parte delle realtà organizzative; tuttavia, è necessario che i leader di settore
continuino a promuoverne lo sviluppo per poter ottenere, quanto prima, sistemi solidi ed efficaci che possano rivoluzionare concretamente
il mondo della sicurezza e contrastare il crescente impiego di intelligenza artificiale offensiva da parte dei cybercriminali.
